% In this file you should put all LaTeX macros and settings to be used both by
% the pdf version and the web version.
% This should be most of your macros.

% The theorem-like environments defined below are those that appear by default
% in the dependency graph. See the README of leanblueprint if you need help to
% customize this. 
% The configuration below use the theorem counter for all those environments
% (this is what the [theorem] arguments mean) and never resets it.
% If you want for instance to number them within chapters then you can add
% [chapter] at the end of the next line.
\newtheorem{theorem}{Theorem}
\newtheorem{proposition}[theorem]{Proposition}
\newtheorem{lemma}[theorem]{Lemma}
\newtheorem{corollary}[theorem]{Corollary}

\theoremstyle{definition}
\newtheorem{definition}[theorem]{Definition}

% Protocol-dynamics commentary: expository blocks that no formal content depends on.
% To suppress every such block throughout the document, comment out the first
% definition and uncomment the second.
\newcommand{\dynamics}[2]{\par\medskip\noindent\textbf{#1.}\ #2\par\medskip}
%\newcommand{\dynamics}[2]{}
% Reference to an algorithm float: the full edition renders the numbered
% reference; the default edition, which omits the floats, renders the prose
% name (arg 2) instead (its roots \renewcommand this).
\newcommand{\algoref}[2]{Algorithm~\ref{#1}}
% The coin's delivery-failure mass, kept typographically apart from the Dirac
% distribution \delta(-). The corresponding Lean field of Params is named δ.
\newcommand{\failmass}{\delta_{\mathrm{f}}}
% The gather-based chain: each of its protocol-shaped systems carries the name
% of its ladder-chain counterpart under a superscript naming the source. The
% Lean development writes that superscript as the namespace AFW.
\newcommand{\afw}[1]{#1^{\mathsf{AFW}}}
\newcommand{\abdy}[1]{#1^{\mathsf{ABDY}}}
